% ---------------------------------------------------------------------------
% Standalone section: a run of the sign-flipped BKNS algorithm that gets stuck.
% Uses only macros already in preamble.tex (\items, \agents, \cost, \subsidy,
% \optsubsidy, \set, \abs, \Z, \Mset). Drop into sections/ and \input it.
% ---------------------------------------------------------------------------

\section{Appendix}
\label{sec:whynotbkns}

% \subsection{The instance}

% Three agents $\agents = \set{1,2,3}$ and three chores $\items = \set{a,b,c}$.
% Each cost function depends only on the size of the bundle:
% \[
%   \cost_1(S) \;=\; \max\bigl(0,\ \abs S - 1\bigr),
%   \qquad
%   \cost_2(S) \;=\; \cost_3(S) \;=\; \min\bigl(\abs S,\,2\bigr).
% \]
% Every marginal is $0$ or $1$ and $\cost_i(\emptyset) = 0$, so the instance is
% negative dichotomous.

% \begin{center}
% \begin{tabular}{c|cccc}
% $\abs S$ & $0$ & $1$ & $2$ & $3$ \\
% \hline
% $\cost_1(S)$ & $0$ & $0$ & $1$ & $2$ \\
% $\cost_2(S) = \cost_3(S)$ & $0$ & $1$ & $2$ & $2$ \\
% \end{tabular}
% \end{center}

% \subsection{The algorithm, sign-flipped}

% Algorithm~1 of~\cite{BKNS22} allocates one item per iteration and maintains an
% envy-free solution with subsidies in $\set{0,1}^n$. Its EXTEND branch requires
% the incoming item to have marginal $+1$ for the recipient. No chore has a
% positive marginal, so under the sign flip that branch never applies and every
% iteration goes to FINDSINK, which assigns the item to an agent of
% \emph{maximum} subsidy, i.e.\ an agent in $\Mset(\subsidy)$.

% \subsection{The run}

% \begin{center}
% \begin{tabular}{llcc}
% step & allocation & $\optsubsidy$ & $\Mset(\optsubsidy)$ \\
% \hline
% start & $(\emptyset,\ \emptyset,\ \emptyset)$ & $(0,0,0)$ & $\set{1,2,3}$ \\
% $a$ to agent $1$ & $(\set{a},\ \emptyset,\ \emptyset)$ & $(0,0,0)$ & $\set{1,2,3}$ \\
% $b$ to agent $1$ & $(\set{a,b},\ \emptyset,\ \emptyset)$ & $(1,0,0)$ & $\set{1}$ \\
% \end{tabular}
% \end{center}

% Both steps are legal: at each one the recipient lies in $\Mset(\optsubsidy)$
% and the resulting allocation is envy-freeable with subsidies in $\set{0,1}^3$.

% \subsection{The last chore}

% At this point $\Mset(\optsubsidy) = \set 1$, so FINDSINK must give $c$ to agent
% $1$. It cannot. Nor can any other agent, and nor does reassigning the bundles
% help:

% \begin{center}
% \begin{tabular}{lll}
% $c$ added to & resulting bundles & least $\optsubsidy$ over all reassignments \\
% \hline
% agent $1$'s bundle & $\set{a,b,c},\ \emptyset,\ \emptyset$ & $(2,0,0)$ \\
% agent $2$'s bundle & $\set{a,b},\ \set{c},\ \emptyset$ & $(2,1,0)$ \\
% agent $3$'s bundle & $\set{a,b},\ \emptyset,\ \set{c}$ & $(2,0,1)$ \\
% \end{tabular}
% \end{center}

% Every option requires a subsidy of $2$. The algorithm is stuck.

% \subsection{The instance itself is fine}

% The allocation $(\set a,\ \set b,\ \set c)$ gives every agent cost $1$ and is
% envy-free with $\subsidy = \mathbf 0$. What fails is the one-item-at-a-time
% construction: reaching it requires taking back the two chores already placed on
% agent $1$.


\begin{example}[Failure of a BKNS-type (\cite{BKNS22}) incremental extension]
\label{app:extension-failure}

Consider three agents $N=\{1,2,3\}$ and three chores $M=\{a,b,c\}$. We work in cost form, writing
\[
c_i=-v_i
\]
for agent \(i\)'s non-negative cost function. Thus, an allocation \(\mathcal A=(A_1,\ldots,A_n)\) together with a subsidy vector \(p=(p_1,\ldots,p_n)\) is envy-free iff
\[
c_i(A_i)-p_i\le c_i(A_j)-p_j
\qquad\text{for all } i,j\in N.
\]
The cost functions depend only on bundle size: \[ c_1(S)=\max\{0,|S|-1\}, \qquad c_2(S)=c_3(S)=\min\{|S|,2\}. \] Thus, \[
\begin{array}{c|cccc}
|S| & 0 & 1 & 2 & 3\\
\hline
c_1(S) & 0 & 0 & 1 & 2\\
c_2(S)=c_3(S) & 0 & 1 & 2 & 2 .
\end{array}
\] Every marginal cost belongs to $\{0,1\}$, so these are dichotomous cost functions. Suppose  $c$ is currently unallocated, and the partial allocation is \[
X_1=\{a,b\},
\qquad
X_2=X_3=\emptyset.
\] Note that this partial allocation is itself reachable from the empty allocation through successive insertions while maintaining the unit-subsidy invariant:
\[
(\emptyset,\emptyset,\emptyset)
\longrightarrow
(\{a\},\emptyset,\emptyset)
\longrightarrow
(\{a,b\},\emptyset,\emptyset),
\]
with pointwise-minimal subsidy vectors
\(
(0,0,0), (0,0,0),\text{ and } (1,0,0),
\) respectively. The pointwise-minimal subsidy vector of $(X_1,X_2,X_3)$ is
\[
p^*=(1,0,0).
\]

Indeed, agent $1$ has cost $1$ for her own bundle and cost $0$ for either empty bundle, so she requires one unit of subsidy. Agents $2$ and $3$ have no envy, since they have cost $0$ for their own empty bundles and cost $2$ for $X_1$.

We claim that the remaining chore $c$ cannot be inserted so that the resulting allocation admits an envy-free subsidy vector in $\{0,1\}^3$, even if the existing bundles $\{a,b\},\emptyset,\emptyset$ may first be permuted among the agents.

If $c$ is added to the two-chore bundle, the resulting bundle sizes are $(3,0,0)$. Every agent evaluates a three-chore bundle at cost $2$ and an empty bundle at cost $0$. Hence, the agent receiving the three-chore bundle requires a subsidy at least two units larger than that of an agent receiving an empty bundle. Thus, no subsidy vector in $\{0,1\}^3$ can make the resulting allocation envy-free.

It remains to consider the case in which $c$ is added to an empty bundle. The resulting bundle sizes are $(2,1,0)$, up to permutation. If an agent in $\{2,3\}$ receives the two-chore bundle, she evaluates it at cost $2$ and the empty bundle at cost $0$, again requiring a subsidy difference of at least two.

Therefore, the only remaining possibility is that agent $1$ receives the two-chore bundle. Let $i\in\{2,3\}$ receive the singleton and let $j$ receive the empty bundle. Envy-freeness for agent $1$ with respect to $i$ requires
\[
1-p_1\le 0-p_i,
\]
and hence
\[
p_1\ge p_i+1.
\]
Envy-freeness for agent $i$ with respect to $j$ requires
\[
1-p_i\le -p_j,
\]
and hence
\[
p_i\ge p_j+1.
\]
Combining the two inequalities gives
\[
p_1\ge p_j+2,
\]
which is impossible for $p\in\{0,1\}^3$. Hence, no assignment of the additional chore $c$, even after an arbitrary permutation of the previously allocated bundles, preserves the unit-subsidy guarantee. Notice that this is an obstruction to the incremental construction, not to the instance itself. The complete allocation
\[
(\{a\},\{b\},\{c\})
\]
is envy-free without subsidies: agent $1$ evaluates every singleton at cost $0$, while agents $2$ and $3$ evaluate every singleton at cost $1$. \hfill $\square$
\end{example}